\par\medskip\noindent
\SGAInlineRefTarget{sga3:xix:lemma:2-3}\textbf{\textit{\SGARefLink{sga3:xix:lemma:2:3}{Lemma 2.3}.}}---
\textit{Let \(S\) be a scheme, let \(G\) be an \(S\)-group scheme that
is smooth and affine over \(S\), with connected fibers, and let \(T\)
be a torus of \(G\).  The set of points \(s\in S\) for which
\(G_{\bar s}\) is reductive, of semisimple rank \(1\), with maximal
torus \(T_{\bar s}\), is an open subset \(U\) of \(S\).}
\SGARefTarget{sga3:xix:lemma:2:3}

\SGAProof{Proof}
Since \(G\) and \(T\) are flat over \(S\), the function
\[
  s
  \longmapsto
  \dim(G_{\bar s}/\bar s)-\dim(T_{\bar s}/\bar s)
\]
is locally constant.  Let \(U_1\) be the open subset of \(S\) on which
this function has the value \(2\).

\footnote{Editor's note: the following details have been added.  Note
also that §6 is independent of the rest of this Exposé.}
By \SGARefLink{sga3:xix:proposition:6-3}{6.3} the Weyl group
\[
  W_G(T)
  =
  \underline{\operatorname{Norm}}_{G}(T)/
  \underline{\operatorname{Centr}}_{G}(T)
\]
is represented by an \(S\)-group scheme that is étale and separated
over \(S\), and the function
\[
  s\longmapsto
  \text{the number of points of \(W_G(T)_{\bar s}\)}
\]
is lower semicontinuous.  Let \(U_2\) be the locus on which this number
is greater than \(1\); it is open.  By Lemma 1.11, the locus on which
\(G_{\bar s}\) is reductive, of semisimple rank \(1\), with maximal
torus \(T_{\bar s}\), is
\[
  U=U_1\cap U_2.
\]
Furthermore, \(W_G(T)_{\bar s}\) has exactly two points for every
\(s\in U\).  Consequently, by SGA 1, I, 10.9 and
EGA IV\textsubscript{3}, 15.5.1 and IV\textsubscript{4}, 18.12.4,
\(W_G(T)_U\) is finite and étale over \(U\).

\par\medskip\noindent
\SGAInlineRefTarget{sga3:xix:remark:2-4}\textbf{\textit{\SGARefLink{sga3:xix:remark:2:4}{Remark 2.4}.}}---
\textit{The group \(W_G(T)_U\) is isomorphic to
\((\mathbb Z/2\mathbb Z)_U\).}
\SGARefTarget{sga3:xix:remark:2:4}

Indeed, the preceding argument shows that it is finite and étale over
\(U\).  Since the automorphism functor of
\((\mathbb Z/2\mathbb Z)_U\) is the trivial group, it is enough to
verify the assertion locally.  It is immediate whenever the algebra
\(\mathscr A\) defining \(W_G(T)_U\) is a free
\(\mathcal O_U\)-module.\footnote{Editor's note: the hypothesis implies
that \(\mathscr A\) is locally free of rank \(2\).  Since the
augmentation ideal \(\mathscr I\) is a direct summand of \(\mathscr A\),
after replacing \(U\) by a sufficiently small affine open
\(\operatorname{Spec}(R)\), one may suppose that
\(I=\Gamma(U,\mathscr I)\) is a free \(R\)-module of rank \(1\).
If \(x\) is a generator of \(I\), then \(x^2=ax\) for some \(a\in R\);
the assumption that \(A=\Gamma(U,\mathscr A)\) is étale implies that
\(a\) is invertible, which gives the asserted local identification.}
